\section{支撑材料与结果生成链}

论文入口为\texttt{main.tex}，正文四至七章分别为四个问题，每章依次包含分析、模型建立和求解。原\texttt{example.tex}、类文件及样式文件保留。计算代码、原始模板和第一版源代码分别保存在\texttt{code/}、\texttt{originals/}、\texttt{versions/v1\_before\_settlement\_revision/}；分问重排前的正文与代码另存\texttt{versions/v2\_before\_question\_restructure/}。修订稿未编译，模板示例PDF及既有论文PDF均不能视为本次源稿的编译结果。

\begin{table}[htbp]\centering\small
\caption{修订稿计算入口与结果文件}\label{tab:files}
\begin{tabularx}{\linewidth}{p{0.43\linewidth}X}\toprule
文件 & 用途\\\midrule
\texttt{code/data\_forecast.py} & 原始附件读取及因果预测基础\\
\texttt{code/optimize\_dispatch.py} & 保留原版LP、残差场景与规则执行\\
\texttt{code/dispatch\_v2.py} & 修订控制器、跨日续接与逐段结算\\
\texttt{code/calibrate\_v2.py} & 一月费用选型、匹配参数与冻结配置\\
\texttt{code/run\_main\_v2.py} & 四组冻结主案及唯一导出数据源\\
\texttt{code/experiments\_v2.py} & 场景、风险、备用、反馈及控制对照\\
\texttt{code/diagnose\_execution\_gap.py} & 原版全决策点规划与规则回放诊断\\
\texttt{code/diagnose\_inner\_v2.py} & 指定四日的同合同两控制器诊断\\
\texttt{code/time\_index.py} & 物理时间、模板列序及分组映射\\
\texttt{code/make\_paper\_assets\_v2.py} & 正文数表、插图及逐日CSV\\
\texttt{code/make\_question\_tables.py} & 分问编排中的预测与场景对照表\\
\texttt{code/make\_visual\_atlas.py} & 热力矩阵、瀑布图与场景阴影带\\
\texttt{code/scenario\_convergence.py} & 真实场景缩减迭代记录与收敛分位带\\
\texttt{code/check\_dispatch\_v2.py} & 数值守恒、未来扰动及价格机制检查\\
\texttt{code/export\_workbooks.mjs} & 导入原模板并填写五本结果\\
\texttt{code/verify\_workbooks.py} & 工作簿逐值、标签、日期、汇总检查\\
\texttt{code/run\_revision.ps1} & 默认复用冻结配置的修订复算入口\\
\bottomrule\end{tabularx}
\end{table}

\texttt{selection\_v2.json}保存每个一月候选、费用、末库存、最终配置及固定参数匹配值。\texttt{prepared\_v2\_*.npz}为已选预测输入，\texttt{*\_v2.npz}为334天执行结果；旧\texttt{*\_full.npz}是原版365天记录，提供冻结的一月运行前缀，不作为修订主结果使用。

每个主结果中，\texttt{plan}为零时合同，\texttt{adjusted}为每段最终有效常规量，\texttt{charge/discharge}为电池外侧实际电量，\texttt{energy}每一天含145个边界状态，\texttt{reference\_energy}单独保存外层参考路径。\texttt{emergency/spill}分别为紧急与弃余电量；\texttt{plan\_cost}、\texttt{adjustment\_cost}及\texttt{emergency\_cost}之和为\texttt{total\_cost}。所有量保持完整精度，表内显示再按相应小数位舍入。

主执行数组汇入\texttt{data/results\_export.json}，其版本标记为\texttt{v2\_settlement\_control}。工作簿与正文数表均使用同一批执行数组和统一索引，不从PDF抄录汇总数值。\texttt{daily\_summary.csv}保存逐日费用和库存；\texttt{experiments\_v2.json/npz}保存配置及逐日对照；两个诊断JSON/CSV分别保留规划、回放及场景概率。

\section{时间映射与附件5填写}

附件5购电页保持原模板首标签“0:10--0:20”和末标签次日“0:00--0:10”。首末标签的原始文本在不同工作簿中略有差别，均逐字保留。以下列出源样本、物理模型与模板显示的对应，不据标签再平移数组：

\begin{center}\small
\begin{tabular}{llll}\toprule
数组索引 & 源样本右端点 & 物理时段 & 保留模板显示\\\midrule
$t=0$ & 00:10 & 00:00--00:10 & 0:10--0:20\\
$t=60$ & 10:10 & 10:00--10:10 & 10:10--10:20\\
$t=143$ & 次日00:00 & 23:50--24:00 & 次日0:00--0:10\\
\bottomrule\end{tabular}
\end{center}

具体原标签、Excel行列与四小时块均冻结于\texttt{data/time\_mapping.json/csv}。问题一为纵向144行，其他购电页为横向144列；物理10:00--10:10从同一$t=60$取值。四小时储能区间依次汇总24段，紧急页将同一天相邻正紧急量合并，跨日不合并，无紧急量日期写“无”和0。源样本按右端点平均功率解释是明示建模假设，不将模板显示差异说成题目已给定的物理事实。

计划页的全天购电量及费用分别为$\sum_tg_t$及$\sum_tp_tg_t$；调整页填$a_t$而非$a_t-g_t$，全天费用填式\eqref{eq:settlement}的全部实际账单。没有调整页的问题二与问题四静态案，其完整账单由逐日CSV和正文给出。五本已填文件保存在\texttt{附件/附件5/}，空白原模板仅保存在\texttt{originals/附件5/}，交付时必须区别二者。

本次写入后逐格核对了288720个计划及调整数值：问题一144个，四本年度计划各48096个，两本年度调整各48096个。全部日期与原标签一致、没有空缺，储能和紧急页也分别核对汇总。文件路径、SHA-256与输入JSON指纹存于\texttt{outputs/workbook\_previews/workbook\_file\_manifest.csv}，完整核对记录在同目录的\texttt{workbook\_verification.json}。这些核对验证文件与执行数组的一致性，不能替代独立复现模型和人工论文检查。

\section{数值可行性与因果信息检查}

四个主案逐段供需平衡最大绝对残差不超过$2.3\times10^{-13}$ kWh，状态递推残差不超过$9.1\times10^{-13}$ kWh，日际连续性残差为0。储电量位于1200--10800 kWh，每段充放电不超过$5000/6$ kWh，未出现正充电与正放电同时发生。逐日费用组成残差小于$9\times10^{-12}$元。

在第201日的第63段及其后改变尚未观测的实际负载、光伏、电价和未来发布预报，并重新生成因果预测，规则执行及跨日LP在此前的合同、充放电、紧急量与SOC前缀均不变。该扰动检查连同执行时的信息截取说明代码未因这些未来真值改变已执行动作，但有限检查不等同于任意输入下的信息正确性证明。明细见\texttt{data/numerical\_checks\_v2.json}。

以下节选显示外层交易与内层动作的职责。完整数值实现及参数保存在源文件中。

\begin{lstlisting}[language=Python,basicstyle=\ttfamily\footnotesize]
for t in range(144):
    if t in allowed_updates:
        contract = outer_plan(legal_forecast, observed_energy,
                              original_day_ahead_plan)
    # Contract remains fixed inside the storage controller.
    action = controller(contract[t:], current_net, current_price,
                        latest_legal_future_forecast, observed_energy)
    execute_first_action(action)

plan_cost = sum(price * original_plan)
adjustment_cost = sum(price * (1.5 * maximum(final_plan-original_plan, 0)
                              + 0.5 * maximum(original_plan-final_plan, 0)))
emergency_cost = sum(5 * price * emergency)
total_cost = plan_cost + adjustment_cost + emergency_cost
\end{lstlisting}

\section{指定日期费用分解}

以下附表将正文指定日期的常规量及三项费用并列，便于与工作簿计划页、调整页及紧急页交叉定位。
\begin{table}[htbp]\centering\small
\caption{问题2指定日期计划、调整与费用分解}
\begin{tabularx}{\linewidth}{lRRRRR}\toprule
日期 & 计划量/kWh & 生效量/kWh & 计划费/元 & 调整费/元 & 紧急费/元\\\midrule
2025-03-20 & 67796.64 & 67796.64 & 41219.41 & 0.00 & 0.00\\
2025-06-21 & 40036.30 & 40036.30 & 23986.38 & 0.00 & 0.00\\
2025-09-23 & 69576.50 & 69576.50 & 43098.28 & 0.00 & 3285.98\\
2025-12-21 & 101008.53 & 101008.53 & 65454.44 & 0.00 & 0.00\\
\bottomrule\end{tabularx}\end{table}

AI工具参与的模型整理、程序实现、数值计算及修订范围见\texttt{AI工具使用详情.tex}。人工审阅、排版调整和编译由使用者后续完成，不提前记录为已完成事项。
